% Generated by build/sync_qft_soln.py from committed sources only.
% Source repository: https://github.com/Luca-Yucheng-Jin/QFT-soln
% Source commit: 07d7fe95fd5cf5b26d38a16ee535a04a925277b2
\section{Renormalization (PSI QFT II, Homework 2)}

\noindent\textit{Question prompts paraphrased from the official
\href{https://psi-online.perimeterinstitute.ca/courses/take/quantum-field-theory-ii-student/pdfs/20776395-module-5-homework}{Perimeter PSI Online sheet};
the equations and notation follow the source.}

\subsection{Renormalization for Massive $\phi^4$ Theory}

Work in Euclidean signature and retain terms through one loop. Define the
renormalized action by
\begin{equation}
    S_R[\phi]
    =S[\phi]+\hbar\Delta_1S[\phi]+O(\hbar^2),
\end{equation}
where
\begin{equation}
    S[\phi]
    =\int d^4x\,
    \left[
        \frac{1}{2}(\partial\phi)^2
        +\frac{m_R^2}{2}\phi^2
        +\frac{g_R}{4!}\phi^4
    \right]
\end{equation}
and
\begin{equation}
    \Delta_1S[\phi]
    =\int d^4x\,
    \left[
        \frac{B_1}{2}\phi^2
        +\frac{C_1}{4!}\phi^4
    \right].
\end{equation}
The cutoff-dependent coefficients $B_1$ and $C_1$ are chosen so that physical
quantities expressed in terms of $m_R$ and $g_R$ remain finite as
$\Lambda\to\infty$.

The one-loop four-point 1PI function is
\begin{align}
    \widehat\Gamma_R^{(4)}(p_i)
    &={}
    g_R
    -\frac{\hbar g_R^2}{2}
    \Bigl[
        I(p_1+p_2,m_R)
        +I(p_1+p_3,m_R)
        +I(p_1+p_4,m_R)
    \Bigr]
    \notag\\
    &\quad+\hbar C_1+O(\hbar^2).
\end{align}
Using a hollow square for the quartic counterterm, these contributions have the
following topologies:
\begin{center}
\begin{tikzpicture}[
    baseline=(current bounding box.center),
    psihwTwoScalar/.style={scalar, semithick},
    psihwTwoVertex/.style={circle, fill=black, draw=black,
                           minimum size=3.4pt, inner sep=0pt},
    psihwTwoFourCT/.style={rectangle, fill=white, draw=black, semithick,
                          minimum size=6pt, inner sep=0pt},
    every node/.style={font=\scriptsize}
]
\begin{feynman}
    % Tree-level quartic vertex.
    \vertex[psihwTwoVertex] (tree) at (0,0) {};
    \vertex (treeUL) at (-0.65,0.55);
    \vertex (treeLL) at (-0.65,-0.55);
    \vertex (treeUR) at (0.65,0.55);
    \vertex (treeLR) at (0.65,-0.55);
    \diagram*{
        (treeUL) -- [psihwTwoScalar] (tree),
        (treeLL) -- [psihwTwoScalar] (tree),
        (tree) -- [psihwTwoScalar] (treeUR),
        (tree) -- [psihwTwoScalar] (treeLR),
    };
    \node at (1.02,0) {$+$};

    % s-channel bubble.
    \vertex[psihwTwoVertex] (sL) at (1.65,0) {};
    \vertex[psihwTwoVertex] (sR) at (2.55,0) {};
    \vertex (sUL) at (1.20,0.58);
    \vertex (sLL) at (1.20,-0.58);
    \vertex (sUR) at (3.00,0.58);
    \vertex (sLR) at (3.00,-0.58);
    \diagram*{
        (sUL) -- [psihwTwoScalar] (sL),
        (sLL) -- [psihwTwoScalar] (sL),
        (sL) -- [psihwTwoScalar, bend left=55] (sR),
        (sL) -- [psihwTwoScalar, bend right=55] (sR),
        (sR) -- [psihwTwoScalar] (sUR),
        (sR) -- [psihwTwoScalar] (sLR),
    };
    \node at (2.10,-0.88) {$s$};
    \node at (3.38,0) {$+$};

    % t-channel routing.
    \vertex[psihwTwoVertex] (tT) at (4.05,0.42) {};
    \vertex[psihwTwoVertex] (tB) at (4.05,-0.42) {};
    \vertex (tUL) at (3.50,0.78);
    \vertex (tUR) at (4.60,0.78);
    \vertex (tLL) at (3.50,-0.78);
    \vertex (tLR) at (4.60,-0.78);
    \diagram*{
        (tUL) -- [psihwTwoScalar] (tT),
        (tUR) -- [psihwTwoScalar] (tT),
        (tT) -- [psihwTwoScalar, bend left=55] (tB),
        (tT) -- [psihwTwoScalar, bend right=55] (tB),
        (tB) -- [psihwTwoScalar] (tLL),
        (tB) -- [psihwTwoScalar] (tLR),
    };
    \node at (4.05,-1.03) {$t$};
    \node at (4.98,0) {$+$};

    % u-channel routing.
    \vertex[psihwTwoVertex] (uT) at (5.65,0.42) {};
    \vertex[psihwTwoVertex] (uB) at (5.65,-0.42) {};
    \vertex (uUL) at (5.10,0.78);
    \vertex (uUR) at (6.20,0.78);
    \vertex (uLL) at (5.10,-0.78);
    \vertex (uLR) at (6.20,-0.78);
    \diagram*{
        (uUL) -- [psihwTwoScalar] (uT),
        (uLR) -- [psihwTwoScalar] (uT),
        (uT) -- [psihwTwoScalar, bend left=55] (uB),
        (uT) -- [psihwTwoScalar, bend right=55] (uB),
        (uB) -- [psihwTwoScalar] (uLL),
        (uB) -- [psihwTwoScalar] (uUR),
    };
    \node at (5.65,-1.03) {$u$};
    \node at (6.58,0) {$+$};

    % Quartic counterterm.
    \vertex[psihwTwoFourCT] (ctFour) at (7.20,0) {};
    \vertex (ctUL) at (6.55,0.55);
    \vertex (ctLL) at (6.55,-0.55);
    \vertex (ctUR) at (7.85,0.55);
    \vertex (ctLR) at (7.85,-0.55);
    \diagram*{
        (ctUL) -- [psihwTwoScalar] (ctFour),
        (ctLL) -- [psihwTwoScalar] (ctFour),
        (ctFour) -- [psihwTwoScalar] (ctUR),
        (ctFour) -- [psihwTwoScalar] (ctLR),
    };
    \node at (7.20,-0.88) {counterterm};
\end{feynman}
\end{tikzpicture}
\end{center}

\begin{enumerate}[label=\alph*)]
    \item In four-dimensional Euclidean spacetime, write the sharp-cutoff
    integral corresponding to
    \begin{equation}
        I(p_1+p_2,m_R;\Lambda),
    \end{equation}
    where $\Lambda$ regulates the ultraviolet divergence.
    \begin{solution}
        \begin{equation}
        \begin{aligned}
            I(p_1+p_2,m_R;\Lambda) &=
        \int \frac{d^4k}{(2\pi)^4}\frac{1}{(k^2+m_R^2)[(k+p_1+p_2)^2+m_R^2]}\\
        &=\frac{4\pi}{(2\pi)^4}\int_0^\pi d\chi\sin^2 \chi\int^\Lambda_0dk \frac{k^3}{(k^2+m_R^2)(k^2+ s + 2k\sqrt{s}\cos\chi +m_R^2)}.
        \end{aligned}
        \end{equation}
    \end{solution}

    \item Set $m_R^2=0$ and calculate the leading large-$\Lambda$ behavior of
    $I(p,0;\Lambda)$.
    \begin{solution}
        \begin{equation}
            \begin{aligned}
                I(p_1+p_2,0;\Lambda) &= \frac{1}{4\pi^3}\int_0^\Lambda dk\, k \int_{0}^\pi d\chi\frac{\sin^2\chi}{k^2 +s +2k\sqrt{s}\cos\chi}\\
                &= \frac{1}{4\pi^3}\int_0^\Lambda dk\, k \frac{\pi}{4k^2s}\left(k^2 +s - \sqrt{(k^2+s)^2 -4k^2s}\right)\\
                &=\frac{1}{16\pi^2s}\int ^\Lambda_0\frac{1}{k}(k^2+s - |k^2-s|)\\
                &=\frac{1}{16\pi^2s}\left(\int^\Lambda_{\sqrt{s}}dk \frac{2s}{k} + \int_0^{\sqrt{s}}dk \frac{2k^2}{k}\right)\\
                &=\frac{1}{16\pi^2s}\left(1 + \log\frac{\Lambda^2}{s}\right).
            \end{aligned}
        \end{equation}
    \end{solution}

    \item By differentiating under the integral, show that
    \begin{equation}
        \frac{\partial}{\partial(m_R^2)}I(p,m_R;\Lambda)
    \end{equation}
    is finite as $\Lambda\to\infty$. Conclude that the divergent part in the
    massive theory is the same as the one identified in the massless theory.
    \begin{solution}
        By differentiating with respect to $m_R^2$, the integral now looks like
        \begin{equation}
            \begin{aligned}
                \sim\int^\Lambda_0 dk \frac{k^3}{k^6},
            \end{aligned}
        \end{equation}
        which is finite. Since the rate of change of the integral remains finite, the divergent part must be the same for massive and massless theory.
    \end{solution}

    For parts d)--h), work in massless $\phi^4$ theory.

    \item Determine the divergent part of $C_1$ required to make
    $\widehat\Gamma_R^{(4)}(p_i)$ finite and impose
    \begin{equation}
        \widehat\Gamma_R^{(4)}(p_i^{\mathrm{ref}})=g_R,
        \qquad
        |p_1^{\mathrm{ref}}+p_2^{\mathrm{ref}}|
        =|p_1^{\mathrm{ref}}+p_3^{\mathrm{ref}}|
        =|p_1^{\mathrm{ref}}+p_4^{\mathrm{ref}}|
        =\mu.
    \end{equation}
    Only the terms in the counterterm that diverge as $\Lambda\to\infty$ need
    be given explicitly.
    \begin{solution}
        Using the result from b), we have
        \begin{equation}
            g_R - \frac{\hbar g_R^2}{2}\frac{3}{16\pi^2}\left(1 + \log{\frac{\Lambda^2}{\mu^2}}\right) + \hbar C_1 = g_R \implies C_1 = \frac{3g_R^2}{32\pi^2}\left(1 + \log{\frac{\Lambda^2}{\mu^2}}\right).
        \end{equation}
    \end{solution}

    \item Find $\widehat\Gamma_R^{(4)}(p_i,g_R,\mu)$ at general external momenta.
    Renormalization-scale independence requires
    \begin{equation}
        \widehat\Gamma_R^{(4)}(p_i,g_R,\mu)
        =\widehat\Gamma_R^{(4)}(p_i,g_R',\mu'),
    \end{equation}
    where the primed parameters are defined at the scale $\mu'$.
    \begin{solution}
        \begin{equation}
        \begin{aligned}
            \widehat\Gamma_R^{(4)}(p_i,g_R,\mu) &= g_R -\frac{\hbar g_R^2}{32\pi^2}\left(\log\frac{\mu^2}{s} + \log\frac{\mu^2}{t} +\log\frac{\mu^2}{u}\right)\\
            &=g_R' -\frac{\hbar g_R'^2}{32\pi^2}\left(\log\frac{\mu'^2}{s} + \log\frac{\mu'^2}{t} +\log\frac{\mu'^2}{u}\right).
        \end{aligned}
        \end{equation}
    \end{solution}

    \item Evaluate the preceding equality at a convenient momentum configuration
    and derive $g_R'(\mu')$ in terms of $g_R(\mu)$, $\mu$, and $\mu'$.
    \begin{solution}
        Take $s=u=t = \mu'^2$, we find
        \begin{equation}
            g_R'(\mu') = g_R(\mu) -\frac{3\hbar g_R^2(\mu)}{32\pi^2}\log\frac{\mu^2}{\mu'^2}.
        \end{equation}
    \end{solution}

    \item Define the effective coupling by
    \begin{equation}
        g_{\mathrm{eff}}(E)
        =\begin{cases}
            g_R(\mu), & E=\mu,\\
            g_R'(\mu'), & E=\mu',
        \end{cases}
    \end{equation}
    and the beta function by
    \begin{equation}
        \beta(g_{\mathrm{eff}})
        \equiv E\frac{d}{dE}g_{\mathrm{eff}}(E).
    \end{equation}
    Compute the one-loop beta function and plot it as a function of
    $g_{\mathrm{eff}}$.
    \begin{solution}
        The definition here is not very conventional. It is equivalent to defining
        \begin{equation}
            g_\mathrm{eff}(E) = \Gamma_R^{(4)}(E,g_R,\mu),
        \end{equation}
        which immediately implies
        \begin{equation}
            g_{\mathrm{eff}}(E)
            =\begin{cases}
                g_R(\mu), & E=\mu,\\
                g_R'(\mu'), & E=\mu'.
            \end{cases}
        \end{equation}
        Then, we have
        \begin{equation}
            \beta(g_\mathrm{eff}) = \frac{3\hbar g_R^2(\mu)}{16\pi^2} = \frac{3\hbar g_\mathrm{eff}^2}{16\pi^2}.
        \end{equation}
    \end{solution}

    \item Explain what this beta function implies for the interaction strength at
    high and low energies.
    \begin{solution}
        Solving the beta function, we have
        \begin{equation}
            \frac{1}{g_\mathrm{eff}(E)} = \frac{1}{g_0} + \frac{3\hbar}{32\pi^2}\log\frac{\Lambda^2}{E^2}.
        \end{equation}
        Therefore, the theory is free at $IR$ and becomes strongly coupled at the UV.
    \end{solution}
\end{enumerate}

The four-point diagrams introduce no additional divergence in the massive theory,
so the same calculation can be repeated there. For mass renormalization, consider
the one-loop two-point 1PI function
\begin{equation}
    \widehat\Gamma_R^{(2)}(p)
    =p^2+m_R^2
    +\frac{\hbar g_R}{2}T(m_R)
    +\hbar B
    +O(\hbar^2).
\end{equation}
Using a filled square for the quadratic counterterm, its topologies are
\begin{center}
\begin{tikzpicture}[
    baseline=(current bounding box.center),
    psihwTwoScalar/.style={scalar, semithick},
    psihwTwoVertex/.style={circle, fill=black, draw=black,
                           minimum size=3.4pt, inner sep=0pt},
    psihwTwoTwoCT/.style={rectangle, fill=black, draw=black,
                         minimum size=7pt, inner sep=0pt},
    every node/.style={font=\scriptsize}
]
\begin{feynman}
    \vertex (treeL) at (0,0);
    \vertex[psihwTwoVertex] (treeV) at (0.85,0) {};
    \vertex (treeR) at (1.70,0);
    \diagram*{
        (treeL) -- [psihwTwoScalar] (treeV),
        (treeV) -- [psihwTwoScalar] (treeR),
    };
    \node at (2.15,0) {$+$};

    \vertex (tadL) at (2.55,0);
    \vertex[psihwTwoVertex] (tadV) at (3.40,0) {};
    \vertex (tadR) at (4.25,0);
    \diagram*{
        (tadL) -- [psihwTwoScalar] (tadV),
        (tadV) -- [psihwTwoScalar] (tadR),
    };
    \draw[psihwTwoScalar]
        (tadV) ++(0,0.30) circle (0.30);
    \node at (4.70,0) {$+$};

    \vertex (ctL) at (5.10,0);
    \vertex[psihwTwoTwoCT] (ctV) at (5.95,0) {};
    \vertex (ctR) at (6.80,0);
    \diagram*{
        (ctL) -- [psihwTwoScalar] (ctV),
        (ctV) -- [psihwTwoScalar] (ctR),
    };
\end{feynman}
\end{tikzpicture}
\end{center}

\begin{enumerate}[label=\alph*), resume]
    \item Write the sharp-cutoff integral $T(m_R;\Lambda)$ that regulates the UV
    divergences.
    \begin{solution}
        \begin{equation}
            T(m_R;\Lambda) = \int \frac{d^4k}{(2\pi)^4}\frac{1}{k^2 +m^2} = \frac{1}{8\pi^2}\int^\Lambda_0dk \frac{k^3}{k^2+m^2.}
        \end{equation}
    \end{solution}

    \item Evaluate $T(m_R;\Lambda)$ at large $\Lambda$. Identify separately the
    quadratic divergence, logarithmic divergence, and term that remains finite as
    $\Lambda\to\infty$.

    \textit{Hint:} Radial coordinates may be useful. The area of the unit
    $n$-sphere is
    \begin{equation}
        \int d^n\Omega
        =\frac{2\pi^{n/2}}{\Gamma(n/2)}.
    \end{equation}
    \begin{solution}
        \begin{equation}
            T(m_R;\Lambda) = \frac{1}{16\pi^2}\left(\Lambda^2 - m_R^2\log \frac{\Lambda^2 +m_R^2}{m_R^2}\right) =  \frac{1}{16\pi^2}\left(\Lambda^2 - m_R^2\log \frac{\Lambda^2 }{m_R^2} + \mathcal{O}(\Lambda^{-2})\right).
        \end{equation}
    \end{solution}

    The physical mass is determined by the pole of the propagator, equivalently
    by
    \begin{equation}
        \widehat\Gamma_R^{(2)}
        (p_i^{\mathrm{ref}},m_R,g_R,\mu)=0,
        \qquad
        |p_i^{\mathrm{ref}}|=m_{\mathrm{phys}},
    \end{equation}
    for any $\mu$. For practical subtraction, impose instead
    \begin{equation}
        \widehat\Gamma_R^{(2)}(0,m_R,g_R,\mu)=m_R^2,
        \qquad
        \mu=m_R,
    \end{equation}
    and take the mass counterterm to have the form
    \begin{equation}
        B(m_R,g_R,\mu,\Lambda)
        =B_0(g_R,\mu,\Lambda)
        +m_R^2B_1(g_R,\mu,\Lambda).
    \end{equation}

    \item Determine $B_0(g_R,\mu,\Lambda)$ and
    $B_1(g_R,\mu,\Lambda)$ so that the preceding renormalization condition holds.
    \begin{solution}
        We can rewrite
        \begin{equation}
            \widehat\Gamma_R^{(2)}(p,m_R,g_R,\mu) = p^2 + m_R^2 + \hbar\left(\frac{g_R(\mu)\Lambda^2}{32\pi^2} + B_0 + m_R^2\left(-\frac{g_R(\mu)}{32\pi^2}\log\frac{\Lambda^2}{\mu^2} - \frac{g_R(\mu)}{32\pi^2}\log\frac{\mu^2}{m_R^2}+B_1\right)\right).
        \end{equation}
        The renormalisation condition implies
        \begin{equation}
             B_1(g_R,\mu,\Lambda) = \frac{g_R}{32\pi^2}\log\frac{\Lambda^2}{\mu^2} , \quad B_0(g_R,\mu,\Lambda) = -\frac{g_R\Lambda^2}{32\pi^2}.
        \end{equation}
    \end{solution}

    \item Check that these counterterms make
    $\widehat\Gamma_R^{(2)}(p_i,m_R,g_R,\mu)$ finite for arbitrary $m_R$ and $p$.
    \begin{solution}
        \begin{equation}
            \widehat\Gamma_R^{(2)}(p,m_R,g_R,\mu) = p^2 + m_R(\mu)^2 -\frac{\hbar g_R(\mu)}{32\pi^2}\log\frac{\mu^2}{m_R(\mu)^2}.
        \end{equation}
    \end{solution}

    \item Require $\widehat\Gamma_R^{(2)}$ to be independent of the
    renormalization scale, as for the four-point function, and derive an equation
    relating $[m_R'(\mu')]^2$ and $[m_R(\mu)]^2$. Retain terms only through
    order $\hbar$ and express the result in terms of $g_R$.
    \begin{solution}
        \begin{equation}
                    m_R^2(\mu)
        -
        \frac{\hbar g_R(\mu)m_R^2(\mu)}{32\pi^2}
        \log\left(\frac{\mu^2}{m_R^2(\mu)}\right)
        =
        [m_R'(\mu')]^2
        -
        \frac{\hbar g_R'(\mu')[m_R'(\mu')]^2}{32\pi^2}
        \log\left(\frac{\mu'^2}{[m_R'(\mu')]^2}\right).
        \end{equation}
                This implies
                \begin{equation}
                            [m_R'(\mu')]^2
=
m_R^2(\mu)
+
\frac{\hbar g_R(\mu)m_R^2(\mu)}{32\pi^2}
\log\left(\frac{\mu'^2}{\mu^2}\right)
+O(\hbar^2).
                \end{equation}
    \end{solution}

    \item Show that when $\mu^2=m_R^2$, the renormalized mass agrees in magnitude
    with the physical mass:
    \begin{equation}
        |m_R|=|m_{\mathrm{phys}}|.
    \end{equation}
    \begin{solution}
        \begin{equation}
            \widehat\Gamma_R^{(2)}(|m_\mathrm{phys}|,m_R = \mu,g_R,\mu) = 0 \implies |m_\mathrm{phys}| = |m_R|.
        \end{equation}
    \end{solution}

    \item Calculate the anomalous dimension of the mass-squared operator,
    \begin{equation}
        \gamma_{m^2}(g_R)
        \equiv
        \mu\frac{\partial}{\partial\mu}\ln m_R^2(\mu).
    \end{equation}
    \begin{solution}
    Starting from the physical mass
    \begin{equation}
        m_{\mathrm{phys}}^2
=
m_R^2(\mu)
\left[
1-
\frac{\hbar g_R(\mu)}{32\pi^2}
\log\left(\frac{\mu^2}{m_R^2(\mu)}\right)
\right]
+O(\hbar^2),
    \end{equation}
    and hence
    \begin{equation}
        \log m_{\mathrm{phys}}^2
=
\log m_R^2(\mu)
-
\frac{\hbar g_R(\mu)}{32\pi^2}
\log\left(\frac{\mu^2}{m_R^2(\mu)}\right)
+O(\hbar^2).
    \end{equation}
    Then,
        \begin{equation}
            0 = \frac{d\log m_\mathrm{phys}^2}{d\log \mu} = \gamma_{m^2}(g_R) - \frac{2\hbar g_R}{32\pi^2} + \mathcal{O}(\hbar^2),
        \end{equation}
        which implies
        \begin{equation}
                \gamma_{m^2}(g_R)
=
\frac{\hbar g_R}{16\pi^2}
+O(\hbar^2).
        \end{equation}
    \end{solution}
\end{enumerate}

At one loop in $\phi^4$ theory, no wavefunction renormalization is required; the
two-point correction is momentum independent at this order.

\subsection{Renormalization Group Flow for the $O(N)$ Model}

Let $\vec\phi=\{\phi_\alpha\}_{\alpha=1}^{N}$ be an $N$-component scalar field
with $O(N)$ symmetry. In $D$-dimensional Euclidean space, take
\begin{equation}
    S_E[\vec\phi]
    =\int d^Dx\,
    \left[
        \frac{1}{2}(\nabla\vec\phi)^2
        +U(\vec\phi^{\,2})
    \right],
\end{equation}
where
\begin{equation}
    X\equiv\vec\phi^{\,2}
    =\sum_{\alpha=1}^{N}\phi_\alpha^2,
\end{equation}
and $U(X)$ is regular. Restrict the Fourier support to $|k|\leq\Lambda$ and
write
\begin{equation}
    Z
    =\int\mathcal D_\Lambda[\vec\phi]\,
    \exp\!\left[-\frac{1}{\hbar}S_E[\vec\phi]\right].
\end{equation}
Separate the field into slow and fast components,
\begin{equation}
    \vec\phi
    =\widetilde{\vec\phi}+\vec\phi_{\mathrm{fast}},
\end{equation}
where $\widetilde{\vec\phi}$ has modes $|k|\leq\Lambda/S$ and
$\vec\phi_{\mathrm{fast}}$ has modes $\Lambda/S<|k|\leq\Lambda$. Integrate a
thin shell using
\begin{equation}
    S=1+ds,
    \qquad
    0<ds\ll1.
\end{equation}

\begin{enumerate}[label=\alph*)]
    \item Hold $\widetilde{\vec\phi}$ fixed and expand $U(X)$ through second
    order in $\vec\phi_{\mathrm{fast}}$ (the Gaussian approximation), expressing
    the result in terms of $U'$ and $U''$. Keep the vector structure explicit.

    \item In the local-potential approximation, take
    $\widetilde{\vec\phi}$ to be constant, so only its $k=0$ mode is nonzero.
    Show that the fast-mode integral separates into one longitudinal mode parallel
    to $\widetilde{\vec\phi}$ and $N-1$ decoupled transverse modes.

    \item Show that the two kinds of mode renormalize the potential according to
    \begin{equation}
        U\longrightarrow U_{\mathrm{eff}}
        =U+ds\left(
            \Delta U_{\parallel}
            +(N-1)\Delta U_{\perp}
        \right),
    \end{equation}
    and calculate $\Delta U_{\parallel}(X)$ and $\Delta U_{\perp}(X)$.

    \item Show that the rescalings
    \begin{equation}
        \vec\phi\longrightarrow
        (1-ds)^{(2-D)/2}\vec\phi,
        \qquad
        x\longrightarrow(1-ds)x
    \end{equation}
    restore the cutoff to $\Lambda$.

    \item Determine the local-potential flow equation
    \begin{equation}
        S\frac{\partial}{\partial S}U_S(X)=?
    \end{equation}
    for the $O(N)$ model. As in the lecture, set $\Lambda=1$ if desired.

    \item Check that $N=1$ reproduces the one-component flow for $V(\phi)$,
    using $U(\phi^2)=V(\phi)$.

    \item Truncate the potential at second order in $X=\vec\phi^{\,2}$,
    \begin{equation}
        U(X)=\frac{t}{2}X+\frac{u}{8}X^2.
    \end{equation}
    Expand the flow equation consistently through $X^2$. Derive both the flow
    equations and the Wilson functions $S\,\partial g_S/\partial S$ for the
    renormalized couplings $t_S$ and $u_S$.
\end{enumerate}
